% GENERATED FILE. Edit canonical lessons, then run npm run generate:books.
% canonical-source-hash: fnv1a64:1428ff501bc8cdd1
% canonical-lessons: RU-C47-sugar, RU-C47-honey, RU-C47-rice, RU-C47-mushroom, RU-C47-glass

\chapter{Sugar, Honey, Rice, Mushrooms and a Glass}
\label{ch:ru-sugar-honey-rice-mushrooms-and-a-glass}
\hlchaptermodality{47}

\begin{chapteropening}
\textbf{By the end of this chapter:} \emph{I can name sugar, honey, rice, a mushroom and a glass.}

Complete the last lesson of chapter 47: I can name sugar, honey, rice, a mushroom and a glass.
\end{chapteropening}

\section[\texorpdfstring{\ru{сахар} (sákhar)}{сахар (sákhar)}]{\ru{сахар} (sákhar) — sugar}

\label{lesson:RU-C47-sugar}



\begin{quote}
Before the new one: say the Russian for porridge, then the Russian for a pie.
\end{quote}

\subsection*{You'll want to know: \ru{сахар}}
\textbf{\ru{сахар}} (\emph{sákhar}) — \textquotedblleft{}sugar\textquotedblright{}.

\begin{cousinweb}
\textbf{\ru{сахар}} $\leftarrow$ Greek \emph{sákcharon}, from Pali \emph{sakkharā}, Sanskrit \emph{śarkarā}, \textquotedblleft{}gravel, grit, sugar\textquotedblright{} (Vasmer). The same Sanskrit word, through Arabic, gave English \emph{sugar}.
\end{cousinweb}

\begin{grammarlens}[title={What you've built}]
One more word for this chapter.
\end{grammarlens}

\subsection*{Your turn}
\begin{itemize}
\raggedright
  \item \emph{Say it:} \emph{sákhar}
  \item \emph{Say it:} \emph{sákhar}, once more
  \item \emph{Say it:} say \emph{sákhar}
  \item \emph{Recall:} say the Russian for porridge, then the Russian for a pie, then say \emph{sákhar} again
\end{itemize}

\begin{tcolorbox}[breakable,colback=teal!4,colframe=teal!35!black,title={Before you move on}]
What does \ru{сахар} mean? (\textquotedblleft{}Sugar\textquotedblright{}.) Say it once more.
\end{tcolorbox}
\section[\texorpdfstring{\ru{мёд} (myod)}{мёд (myod)}]{\ru{мёд} (myod) — honey}

\label{lesson:RU-C47-honey}



\begin{quote}
Before the new one: say the Russian for a pie, then the Russian for sugar.
\end{quote}

\subsection*{You'll want to know: \ru{мёд}}
\textbf{\ru{мёд}} (\emph{myod}) — \textquotedblleft{}honey\textquotedblright{}.

\begin{cousinweb}
\textbf{\ru{мёд}} $\leftarrow$ Proto-Slavic \emph{*medŭ}, \textquotedblleft{}honey; mead\textquotedblright{} — inherited. Vasmer lists Lithuanian \emph{medùs}, \textquotedblleft{}honey\textquotedblright{}, Sanskrit \emph{mádhu}, \textquotedblleft{}honey, mead\textquotedblright{}, and Greek \emph{méthy}, \textquotedblleft{}wine\textquotedblright{}: English \emph{mead} is a cousin.
\end{cousinweb}

\begin{grammarlens}[title={What you've built}]
One more word for this chapter.
\end{grammarlens}

\subsection*{Your turn}
\begin{itemize}
\raggedright
  \item \emph{Say it:} \emph{myod}
  \item \emph{Say it:} \emph{myod}, once more
  \item \emph{Say it:} say \emph{myod}
  \item \emph{Recall:} say the Russian for a pie, then the Russian for sugar, then say \emph{myod} again
\end{itemize}

\begin{tcolorbox}[breakable,colback=teal!4,colframe=teal!35!black,title={Before you move on}]
What does \ru{мёд} mean? (\textquotedblleft{}Honey\textquotedblright{}.) Say it once more.
\end{tcolorbox}
\section[\texorpdfstring{\ru{рис} (ris)}{рис (ris)}]{\ru{рис} (ris) — rice}

\label{lesson:RU-C47-rice}



\begin{quote}
Before the new one: say the Russian for sugar, then the Russian for honey.
\end{quote}

\subsection*{You'll want to know: \ru{рис}}
\textbf{\ru{рис}} (\emph{ris}) — \textquotedblleft{}rice\textquotedblright{}.

\begin{cousinweb}
\textbf{\ru{рис}} probably came through Low German \emph{rîs} or Dutch \emph{rijs}, from Italian \emph{riso}, from Latin \emph{oryza} and Greek \emph{óryza}, which came out of India (Vasmer).
\end{cousinweb}

\begin{grammarlens}[title={What you've built}]
One more word for this chapter.
\end{grammarlens}

\subsection*{Your turn}
\begin{itemize}
\raggedright
  \item \emph{Say it:} \emph{ris}
  \item \emph{Say it:} \emph{ris}, once more
  \item \emph{Say it:} say \emph{ris}
  \item \emph{Recall:} say the Russian for sugar, then the Russian for honey, then say \emph{ris} again
\end{itemize}

\begin{tcolorbox}[breakable,colback=teal!4,colframe=teal!35!black,title={Before you move on}]
What does \ru{рис} mean? (\textquotedblleft{}Rice\textquotedblright{}.) Say it once more.
\end{tcolorbox}
\section[\texorpdfstring{\ru{гриб} (grib)}{гриб (grib)}]{\ru{гриб} (grib) — a mushroom}

\label{lesson:RU-C47-mushroom}



\begin{quote}
Before the new one: say the Russian for honey, then the Russian for rice.
\end{quote}

\subsection*{You'll want to know: \ru{гриб}}
\textbf{\ru{гриб}} (\emph{grib}) — \textquotedblleft{}a mushroom\textquotedblright{}.

\begin{cousinweb}
\textbf{\ru{гриб}} $\leftarrow$ Proto-Slavic \emph{*gribŭ} — inherited, but its origin is unclear. Vasmer reports a possible link with Lithuanian \emph{griẽti}, \textquotedblleft{}to skim cream\textquotedblright{}, from a sense \textquotedblleft{}slime\textquotedblright{}, among other proposals.
\end{cousinweb}

\begin{grammarlens}[title={What you've built}]
One more word for this chapter.
\end{grammarlens}

\subsection*{Your turn}
\begin{itemize}
\raggedright
  \item \emph{Say it:} \emph{grib}
  \item \emph{Say it:} \emph{grib}, once more
  \item \emph{Say it:} say \emph{grib}
  \item \emph{Recall:} say the Russian for honey, then the Russian for rice, then say \emph{grib} again
\end{itemize}

\begin{tcolorbox}[breakable,colback=teal!4,colframe=teal!35!black,title={Before you move on}]
What does \ru{гриб} mean? (\textquotedblleft{}A mushroom\textquotedblright{}.) Say it once more.
\end{tcolorbox}
\section[\texorpdfstring{\ru{стакан} (stakán)}{стакан (stakán)}]{\ru{стакан} (stakán) — a glass (for drinking)}

\label{lesson:RU-C47-glass}



\begin{quote}
Before the new one: say the Russian for rice, then the Russian for a mushroom.
\end{quote}

\subsection*{You'll want to know: \ru{стакан}}
\textbf{\ru{стакан}} (\emph{stakán}) — \textquotedblleft{}a glass (for drinking)\textquotedblright{}.

\begin{cousinweb}
\textbf{\ru{стакан}} — Old Russian \emph{dostokanŭ}, recorded from the 1300s. It is usually taken from Turkic: Vasmer cites Chagatai \emph{tostakan}, \textquotedblleft{}a wooden bowl\textquotedblright{}, and Kazakh \emph{tustaghan}, \textquotedblleft{}a cup, a ladle\textquotedblright{}.
\end{cousinweb}

\begin{grammarlens}[title={What you've built}]
One more word for this chapter.
\end{grammarlens}

\subsection*{Your turn}
\begin{itemize}
\raggedright
  \item \emph{Say it:} \emph{stakán}
  \item \emph{Say it:} \emph{stakán}, once more
  \item \emph{Say it:} say \emph{stakán}
  \item \emph{Recall:} say the Russian for rice, then the Russian for a mushroom, then say \emph{stakán} again
\end{itemize}

\begin{tcolorbox}[breakable,colback=teal!4,colframe=teal!35!black,title={Before you move on}]
What does \ru{стакан} mean? (\textquotedblleft{}A glass (for drinking)\textquotedblright{}.) Say it once more.
\end{tcolorbox}
